\section{Background}
\label{sec:background}

\subsection{Lloyd's Algorithm and Optimal Quantisation}

Lloyd's algorithm was introduced in the context of pulse-code modulation
(PCM) signal quantisation~\citep{lloyd1982}.
The problem is: given a probability distribution $\rho$ on $\mathbb{R}^d$
and a budget of $k$ codewords, place $k$ points $\{c_1,\ldots,c_k\}$ to
minimise the mean squared quantisation error
\[
  \mathcal{E}(\{c_j\}) = \int_{\mathbb{R}^d}
  \min_{j=1}^{k} \|x - c_j\|^2 \, \rho(x) \, dx.
\]
The optimal solution requires two conditions to hold simultaneously:
(1) each codeword $c_j$ is the centroid (population-weighted mean) of its
Voronoi cell $V_j = \{x : \|x - c_j\| \le \|x - c_i\| \; \forall i\}$;
(2) each Voronoi cell is the set of points nearest to its codeword.
A tessellation satisfying both conditions is a \emph{Centroidal Voronoi
Tessellation} (CVT).

Lloyd's algorithm achieves a CVT iteratively:
\begin{enumerate}
\item Assign each point $x$ to its nearest codeword $\hat{j}(x) = \arg\min_j \|x - c_j\|$.
\item Update each codeword to the centroid of its cell:
  $c_j \leftarrow \int_{V_j} x \, \rho(x) \, dx \,/\, \int_{V_j} \rho(x) \, dx$.
\item Repeat until $\max_j \|c_j^{\mathrm{new}} - c_j^{\mathrm{old}}\| < \epsilon$.
\end{enumerate}
Convergence is guaranteed under mild conditions: Du, Faber, and
Gunzburger~\citeyear{du1999} proved that the energy $\mathcal{E}$ is
non-increasing at each step, and that fixed points are precisely the CVTs.
The algorithm does not guarantee convergence to the global minimum (multiple
local CVTs may exist), but in practice converges quickly for $k \ll n$.

\subsection{Centroidal Voronoi Tessellations: Theory}

Du et al.~\citeyear{du1999} surveyed CVTs comprehensively, establishing:
\begin{itemize}
\item \textbf{Existence}: CVTs exist for any compactly supported $\rho$ and
  any $k \ge 1$.
\item \textbf{Local optimality}: Every CVT is a local minimiser of
  $\mathcal{E}$ in the space of $k$-point configurations.
\item \textbf{Convergence rate}: For smooth $\rho$ on convex domains,
  Lloyd's algorithm converges at a rate $O(1/t)$ in $\mathcal{E}$.
\item \textbf{Applications}: CVTs arise naturally in mesh generation,
  image halftoning, data compression, and spatial clustering.
\end{itemize}

In the redistricting context, $\rho$ is the census population distribution
(supported on census tracts rather than a continuum), the ``points'' are
census tracts, and the ``codewords'' are district seeds.
The geometric compactness of CVT cells --- roughly circular for smooth $\rho$
--- is the primary motivation for \cvd{}: it provides districts that are
compact by construction rather than as a side effect of edge-cut minimisation.

\subsection{K-Medoids Clustering vs.\ K-Means}

The classical k-means algorithm~\citep{lloyd1982} uses the centroid
(population-weighted mean location) of each cluster as the new centre.
In continuous Euclidean space, the centroid is always well-defined.
In discrete graphs, the ``centroid'' of a cluster would be a fractional
position --- a non-existent point with no corresponding tract.

K-medoids clustering resolves this by requiring each centre to be an actual
data point~\citep{park2009}.
The \emph{medoid} of a set $S$ is
\[
  \mathrm{medoid}(S) = \arg\min_{i \in S} \sum_{j \in S} d(i,j),
\]
where $d(\cdot,\cdot)$ is the distance metric.
Unlike the mean, the medoid is always an element of $S$, making it the
natural centre representative for discrete domains.

\cvd{} uses the graph-distance medoid exclusively:
the ``centroid'' of each district after Voronoi assignment is the tract
minimising the sum of BFS hop-count distances to all other tracts in the
district.
This ensures that seeds are always real census tracts, and that each
subsequent Voronoi assignment step starts from a valid, connectable
graph node.

\paragraph{Approximate medoid.}
Computing the exact medoid requires evaluating $\sum_{j \in S} d(i,j)$
for every tract $i \in S$, which costs $O(|S|^2)$ BFS evaluations per
district per iteration.
For large states (NC: $\sim 2{,}660$ tracts; TX: $\sim 5{,}300$ tracts),
the exact medoid is expensive.
\cvd{} uses an approximate medoid: sample up to 50 tracts uniformly at
random from $S$ and compute the BFS sum only to those 50 tracts.
The tract minimising this approximate sum is the approximate medoid.
For $|S| \le 50$, the approximation degenerates to the exact medoid.
For $|S| \gg 50$, the approximation error is bounded by standard
concentration inequalities: the expected error is $O(\sqrt{1/50})$ relative
to the range of per-tract distances.

\subsection{Graph Distance as Redistricting Metric}

BFS hop count on the census-tract adjacency graph is a natural graph-distance
metric for Phase~1 \cvd{}.
Unlike geographic Euclidean distance, BFS hop count:
\begin{itemize}
\item requires no coordinate data (the adjacency graph is already loaded);
\item respects geographic connectivity (two tracts separated by a river
  are farther apart in hop count than two tracts sharing a long boundary,
  even if they are close in Euclidean space);
\item is invariant to state shape and projection (Alaska and Hawaii pose
  no special challenges).
\end{itemize}
The primary limitation of BFS hop count is its insensitivity to geographic
scale: two adjacent rural tracts (large area) are distance 1, as are two
adjacent urban tracts (small area).
Phase~2 geographic \cvd{} (Section~\ref{sec:phase2}) addresses this by
replacing BFS distance with Euclidean distance on projected tract centroids.

\subsection{Structure-Layer Algorithms in the B-Series}

The B-series redistricting platform uses a three-layer compositor:
\emph{Structure} (how each bisection node is solved),
\emph{WeightsOverride} (what edge weights encode), and
\emph{SeedCompositor} (how many seeds are tried and how the best is selected).
\cvd{} is a Structure-layer algorithm.

Table~\ref{tab:structure-comparison-bg} summarises the Structure-layer
algorithms in the B-series and their primary design dimension.

\begin{table}[t]
\centering
\caption{Structure-layer algorithms in the B-series redistricting platform.
  ``Optimises'' describes the primary objective at each bisection node.
  CVD optimises geographic proximity (not edge cut).
  All algorithms enforce population balance via post-hoc rebalance.}
\label{tab:structure-comparison-bg}
\renewcommand{\arraystretch}{1.25}
\begin{tabular}{@{}lllll@{}}
\toprule
Algorithm & Optimises & Randomness & Data needed & Phase \\
\midrule
\metis{}         & Edge cut (graph)       & Seeded  & None         & Done \\
\sa{} (B.19)     & Edge cut (graph) + SA  & Seeded  & None         & Done \\
\geosec{} (B.8)  & Direction + EC ratio   & Seeded  & None         & Done \\
\areasec{} (B.9) & Area-swing + EC        & Seeded  & None         & Done \\
\textbf{\cvd{} (B.22)} & \textbf{Geographic distance} & \textbf{Seeded}
                        & \textbf{None (P1) / Centroids (P2)} & \textbf{This paper} \\
\bfsgrowth{} (B.23) & Population balance + proximity & Seeded & None & B.23 \\
\bottomrule
\end{tabular}
\end{table}

\cvd{} is orthogonal to SeedCompositor strategies: \be{} or \cs{} can be
applied over \cvd{} subproblems just as over \metis{} subproblems.
Running \be{}$(T=50)$ with \cvd{} bisection would sample 50 independent
\cvd{} runs per node (each with a different BFS-derived seed) and select
the geometrically most compact result.
